\chapter{Kết luận}

\section{Tổng kết kết quả đạt được}

Sau quá trình thực hiện, đề tài đã hoàn thành toàn bộ các mục tiêu đề ra ban đầu, cả về mặt kỹ thuật lẫn sản phẩm phần mềm thực tế.

Về mặt kỹ thuật, pipeline đa tác nhân gồm 5 Agent đã hoạt động đúng theo thiết kế đề xuất. Cụ thể, ảnh lá lúa đầu vào trước tiên được đưa qua bước tiền xử lý nhằm chuẩn hóa kích thước và chất lượng hình ảnh, sau đó được phân loại bệnh bằng mô hình Vision Transformer (ViT). Tiếp theo, Agent phân tích hình thái sử dụng Gemini để mô tả các đặc điểm trực quan của vùng bệnh trên lá. Kết quả phân loại được đối chiếu và bổ sung thông tin bằng cách tra cứu tri thức chuyên ngành từ cơ sở dữ liệu vector ChromaDB. Toàn bộ thông tin từ các Agent được tổng hợp lại bởi Agent cuối cùng, tạo ra một báo cáo chẩn đoán hoàn chỉnh bằng tiếng Việt, phù hợp với nhu cầu sử dụng thực tế của người dùng trong nước.

Về hiệu năng mô hình, ViT đạt F1-score trung bình **0.963** trên tập kiểm thử gồm 400 ảnh — một kết quả đáng ghi nhận khi so sánh với mức 0.90–0.92 thường thấy ở các mô hình CNN trong các nghiên cứu tương tự về phân loại bệnh lúa. Tuy nhiên, cần thận trọng khi diễn giải sự chênh lệch này, bởi các nghiên cứu khác nhau thường sử dụng tập dữ liệu và điều kiện đánh giá không hoàn toàn đồng nhất, do đó kết quả không hoàn toàn có thể so sánh trực tiếp.

Về phía sản phẩm phần mềm, hệ thống được xây dựng với định hướng dễ triển khai và dễ sử dụng. Giao diện web được thiết kế bằng tiếng Việt, thân thiện với người dùng không có chuyên môn kỹ thuật. Phần backend cung cấp API RESTful kèm tài liệu được tạo tự động, giúp các nhà phát triển có thể tích hợp hệ thống một cách thuận tiện. Toàn bộ ứng dụng được đóng gói bằng Docker, đi kèm hai script khởi động riêng biệt: `start.sh` dành cho môi trường Linux/macOS và `start.bat` dành cho Windows, giúp việc cài đặt trở nên đơn giản trên nhiều nền tảng khác nhau. Ngoài ra, bộ kiểm thử viết bằng pytest được xây dựng để bao phủ các chức năng cốt lõi của từng Agent, đảm bảo tính ổn định và khả năng kiểm soát chất lượng của hệ thống.

Nhìn lại quá trình thực hiện, thách thức lớn nhất trong dự án không đến từ phía thuật toán hay mô hình học máy, mà chủ yếu xuất phát từ bài toán tích hợp hệ thống. Cụ thể, nhóm phải đảm bảo các Agent giao tiếp với nhau đúng định dạng dữ liệu qua từng bước trong pipeline; đồng thời phải xây dựng cơ chế xử lý lỗi linh hoạt cho những trường hợp Gemini API trả về kết quả không đầy đủ hoặc không đúng cấu trúc mong đợi. Bên cạnh đó, việc tối ưu thứ tự khởi động các container trong môi trường Docker — đảm bảo các dịch vụ phụ thuộc lẫn nhau được khởi động đúng trình tự — cũng là một vấn đề kỹ thuật tốn nhiều công sức giải quyết trong giai đoạn triển khai.
\section{Hạn chế của hệ thống}

Để có cái nhìn toàn diện và trung thực về hệ thống hiện tại, cần thẳng thắn nhìn nhận một số hạn chế còn tồn tại — không phải để phủ nhận giá trị của những gì đã xây dựng, mà để xác định rõ phạm vi áp dụng thực tế và định hướng cải tiến trong tương lai.


Hạn chế lớn nhất và cũng là hạn chế mang tính kiến trúc sâu sắc nhất của hệ thống hiện tại là sự phụ thuộc hoàn toàn vào kết nối Internet. Cụ thể, cả hai thành phần cốt lõi trong pipeline xử lý — Morphology Agent (chịu trách nhiệm phân tích hình thái học của ảnh bệnh lá) và Coordinator Agent (chịu trách nhiệm điều phối luồng xử lý, tổng hợp kết quả và sinh phản hồi ngôn ngữ tự nhiên) — đều phụ thuộc vào việc gọi Gemini API thông qua mạng Internet trong mỗi lần chẩn đoán. Điều này có nghĩa là bất kỳ môi trường nào thiếu kết nối 4G hoặc WiFi ổn định đều khiến toàn bộ hệ thống không thể hoạt động, kể cả khi thiết bị đầu cuối (điện thoại, máy tính bảng) vẫn còn pin và hoàn toàn bình thường. Hệ thống không có cơ chế dự phòng ngoại tuyến (offline fallback), không có khả năng lưu trữ yêu cầu để xử lý sau (queue-and-sync), và không có mô hình cục bộ nào được nhúng sẵn vào ứng dụng. Đây là điểm mâu thuẫn căn bản với nhóm đối tượng người dùng mà hệ thống hướng tới ban đầu — cụ thể là nông dân tại các vùng sâu, vùng xa, miền núi hoặc hải đảo, nơi hạ tầng viễn thông còn yếu kém, tín hiệu mạng chập chờn hoặc hoàn toàn vắng bóng. Chính những nơi đó lại là nơi người nông dân cần được hỗ trợ chẩn đoán bệnh cây nhiều nhất, do thiếu chuyên gia nông nghiệp tại chỗ và khó tiếp cận các dịch vụ tư vấn truyền thống. Vì vậy, đây không đơn thuần là một lỗi kỹ thuật hay một tính năng còn thiếu, mà là một vấn đề kiến trúc căn bản cần được tái thiết kế trong các phiên bản tiếp theo. Các hướng giải quyết tiềm năng bao gồm: triển khai mô hình ngôn ngữ nhỏ gọn (small LLM) trực tiếp trên thiết bị (on-device inference), xây dựng chế độ offline với tập tri thức tĩnh, hoặc áp dụng kiến trúc đồng bộ hóa theo lô (batch sync) khi thiết bị có kết nối trở lại.


Một hạn chế thực tiễn khác cần được nhìn nhận thẳng thắn là độ trễ trong quá trình chẩn đoán. Hiện tại, mỗi lần người dùng gửi ảnh lá cây để chẩn đoán, hệ thống mất khoảng ≈ 15 giây từ lúc nhận ảnh đến lúc trả về kết quả đầy đủ. Khi phân tích chi tiết nguồn gốc của độ trễ này, có thể thấy rõ hai lần gọi Gemini API là nguyên nhân chiếm áp đảo, với tổng thời gian khoảng ≈ 10,6 giây trong tổng số ≈ 14 giây thời gian xử lý thực tế (phần còn lại đến từ các bước tiền xử lý ảnh, truy vấn RAG và render kết quả). Điều này có nghĩa là phần lớn thời gian chờ không phải do năng lực tính toán của thiết bị hay logic nghiệp vụ, mà đến từ độ trễ mạng và thời gian suy luận phía server của Gemini. Ở mức độ sử dụng đơn lẻ — tức là một người dùng chụp một ảnh, chờ kết quả rồi đọc và phản hồi — độ trễ 15 giây vẫn nằm trong ngưỡng có thể chấp nhận được Tuy nhiên, kịch bản này trở nên không phù hợp trong các tình huống thực tế phổ biến hơn: ví dụ, một cán bộ khuyến nông cần chụp và chẩn đoán hàng chục mẫu lá trong một buổi thăm đồng, hay một hệ thống giám sát tự động cần xử lý ảnh liên tiếp từ nhiều camera. Trong những trường hợp đó, độ trễ tích lũy sẽ nhanh chóng trở thành nút thắt cổ chai (bottleneck) làm giảm đáng kể trải nghiệm và hiệu quả sử dụng. Hướng cải thiện rõ ràng nhất là giảm số lần gọi API, ví dụ bằng cách hợp nhất hai lần gọi thành một lần duy nhất thông qua thiết kế prompt thông minh hơn, hoặc sử dụng mô hình LLM chạy cục bộ (local inference) để loại bỏ hoàn toàn độ trễ mạng. Các kỹ thuật như streaming response, kết quả từng phần (partial output), hay hiển thị tiến trình trung gian cũng có thể cải thiện đáng kể cảm nhận của người dùng ngay cả khi tổng thời gian xử lý không đổi.

Hệ thống hiện tại chỉ có khả năng nhận diện và chẩn đoán 4 loại bệnh trên lá lúa. Tuy nhiên, điều quan trọng cần làm rõ là đây không phải là giới hạn của kiến trúc hệ thống — pipeline xử lý hoàn toàn có thể mở rộng để hỗ trợ nhiều loại bệnh hơn — mà là giới hạn đến từ tập dữ liệu huấn luyện của mô hình phân loại ảnh ViT (Vision Transformer) được sử dụng, vốn được lấy từ Hugging Face và chỉ bao gồm 4 nhãn bệnh. Khi có tập dữ liệu huấn luyện phong phú hơn hoặc khi fine-tune lại mô hình với nhiều nhãn hơn, hệ thống hoàn toàn có thể mở rộng phạm vi chẩn đoán mà không cần thay đổi kiến trúc tổng thể. Song song với đó, cơ sở tri thức phục vụ cơ chế RAG (Retrieval-Augmented Generation) hiện tại cũng còn rất hạn chế: chỉ bao gồm 4 tài liệu ngắn, đủ để minh họa và kiểm chứng rằng cơ chế RAG hoạt động đúng về mặt kỹ thuật — tức là hệ thống có thể truy xuất đoạn văn liên quan và đưa vào ngữ cảnh cho LLM — nhưng chưa đủ chiều sâu và chiều rộng để phục vụ một ứng dụng thực tiễn. Một hệ thống tư vấn bệnh cây thực sự cần được hỗ trợ bởi một kho tri thức phong phú hơn nhiều: bao gồm các tài liệu kỹ thuật từ viện nghiên cứu nông nghiệp, hướng dẫn sử dụng thuốc bảo vệ thực vật, kinh nghiệm thực địa được chuẩn hóa, và các tài liệu được cập nhật theo mùa vụ, theo vùng khí hậu. Việc mở rộng cơ sở tri thức RAG là một trong những bước có chi phí thấp nhưng tác động cao nhất để nâng cao chất lượng phản hồi của hệ thống trong các phiên bản tiếp theo.
\section{Hướng phát triển trong tương lai}
Trong số tất cả các hướng cải tiến có thể thực hiện, giải quyết bài toán hoạt động ngoại tuyến (offline) là hướng cấp thiết và có tác động thực tiễn lớn nhất trong giai đoạn trước mắt. Lý do xuất phát trực tiếp từ mâu thuẫn kiến trúc đã được chỉ ra ở phần hạn chế: hệ thống được xây dựng để phục vụ nông dân vùng sâu vùng xa, nhưng lại hoàn toàn phụ thuộc vào kết nối Internet — một nghịch lý cần được tháo gỡ trước khi có thể nói đến việc triển khai thực tế ở quy mô rộng. Một phương án kỹ thuật khả thi và có cơ sở thực tiễn rõ ràng là kết hợp hai mô hình nhỏ gọn chạy hoàn toàn trên thiết bị: mô hình thị giác ViT-Tiny và một mô hình ngôn ngữ nhỏ gọn như Phi-3-mini hoặc Gemma-2B. ViT-Tiny là phiên bản thu gọn của kiến trúc Vision Transformer, với chỉ khoảng 5,7 triệu tham số — nhỏ hơn khoảng 15 lần so với ViT-Base đang được sử dụng trong hệ thống hiện tại. Sự thu gọn về kích thước này giúp mô hình có thể được nạp và chạy suy luận (inference) trực tiếp trên bộ nhớ RAM hạn chế của thiết bị di động mà không gây hiện tượng treo máy hay tràn bộ nhớ, đồng thời vẫn duy trì độ chính xác phân loại ở mức chấp nhận được cho các tác vụ nhận diện bệnh lá cây với số lượng nhãn giới hạn. Về phía mô hình ngôn ngữ, cả Phi-3-mini (phát triển bởi Microsoft) lẫn Gemma-2B (phát triển bởi Google) đều là các mô hình thuộc thế hệ SLM (Small Language Model) được tối ưu hóa đặc biệt cho môi trường tài nguyên thấp. Cả hai đã được kiểm nghiệm thực tế và có thể chạy trơn tru trên các thiết bị điện thoại Android tầm trung — tức là những dòng máy phổ thông có giá từ 3–6 triệu đồng, vốn là thiết bị mà nhiều hộ nông dân hiện đã sở hữu — mà không cần bất kỳ kết nối mạng nào trong quá trình suy luận. Điều này cho phép toàn bộ pipeline từ phân tích ảnh đến sinh phản hồi ngôn ngữ tự nhiên diễn ra hoàn toàn cục bộ trên thiết bị, xóa bỏ hoàn toàn rào cản hạ tầng viễn thông. Ngoài việc giải quyết vấn đề kết nối, phương án này còn mang lại lợi ích phụ đáng kể: giảm độ trễ xuống dưới 5 giây (do không còn độ trễ mạng), bảo vệ dữ liệu người dùng tốt hơn (ảnh không cần gửi lên server bên ngoài), và giảm chi phí vận hành dài hạn khi không phụ thuộc vào API có tính phí theo lượt gọi.


Một trong những điểm mạnh thiết kế rõ ràng nhất của hệ thống hiện tại là khả năng mở rộng phạm vi bệnh mà không cần viết lại bất kỳ đoạn code nào liên quan đến logic Agent. Đây là kết quả trực tiếp của việc áp dụng kiến trúc đa tác nhân (multi-agent) với sự tách biệt rõ ràng giữa các tầng chức năng. Cụ thể, để thêm một loại bệnh mới vào hệ thống, người phát triển chỉ cần thực hiện hai thao tác độc lập: Thứ nhất, thay thế hoặc fine-tune lại mô hình ViT với tập dữ liệu mới bổ sung nhãn bệnh cần thêm. Morphology Agent không quan tâm đến việc mô hình bên dưới nhận diện bao nhiêu loại bệnh — nó chỉ nhận đầu ra phân loại và chuyển tiếp thông tin đó trong pipeline. Do vậy, toàn bộ logic điều phối của hệ thống không bị ảnh hưởng. Thứ hai, bổ sung tài liệu tri thức tương ứng vào ChromaDB — cơ sở dữ liệu vector đang phục vụ cơ chế RAG. Khi có thêm tài liệu mô tả triệu chứng, nguyên nhân và phương pháp xử lý cho loại bệnh mới, Coordinator Agent sẽ tự động truy xuất và tích hợp thông tin đó vào phản hồi mà không cần bất kỳ sự can thiệp lập trình nào. Tính mở rộng này không phải ngẫu nhiên mà có — nó là kết quả của một quyết định kiến trúc có chủ đích: tách biệt hoàn toàn tầng nhận thức thị giác (vision layer), tầng tri thức (knowledge layer) và tầng suy luận ngôn ngữ (reasoning layer). Đây là ưu điểm thiết kế rõ ràng nhất của kiến trúc đa tác nhân, và cũng là lý do để giữ nguyên kiến trúc tổng thể thay vì tái thiết kế từ đầu khi mở rộng hệ thống.

Nhìn xa hơn về lộ trình phát triển, một hướng có giá trị thực tiễn rất cao là tích hợp dữ liệu GPS từ các báo cáo chẩn đoán của người dùng để xây dựng một hệ thống bản đồ dịch tễ bệnh cây trồng theo thời gian thực.
Ý tưởng cốt lõi là: mỗi lần người dùng chụp ảnh và nhận kết quả chẩn đoán, hệ thống — với sự đồng ý của người dùng — ghi lại tọa độ GPS kèm theo loại bệnh được phát hiện và thời điểm xảy ra. Khi dữ liệu được tổng hợp từ hàng trăm, hàng nghìn người dùng trên một vùng địa lý, hệ thống có thể tự động nhận diện các cụm dịch bệnh đang hình thành, theo dõi tốc độ và hướng lan rộng của dịch, và cung cấp cảnh báo sớm đến các hộ nông dân ở vùng lân cận hoặc đến cơ quan quản lý nông nghiệp địa phương. Đây là loại thông tin mà các hệ thống giám sát truyền thống hiếm khi có được do chi phí khảo sát thực địa quá lớn. Tuy nhiên, để hiện thực hóa hướng này, cần giải quyết nghiêm túc ít nhất hai nhóm vấn đề phức tạp. Nhóm thứ nhất là quyền riêng tư dữ liệu. Dữ liệu GPS gắn với thói quen canh tác của từng hộ nông dân là thông tin nhạy cảm, có thể tiết lộ diện tích đất, vị trí thửa ruộng và chu kỳ sản xuất. Hệ thống cần có cơ chế thu thập, có sự đồng thuận rõ ràng, chính sách lưu trữ và ẩn danh hóa dữ liệu minh bạch, và tuân thủ các quy định pháp lý liên quan đến bảo vệ dữ liệu cá nhân. Nhóm thứ hai là độ tin cậy của báo cáo người dùng. Không phải mọi chẩn đoán từ người dùng đều chính xác — ảnh chụp thiếu sáng, góc chụp không chuẩn, hoặc lá bị nhiễm nhiều loại bệnh cùng lúc đều có thể dẫn đến kết quả phân loại sai. Nếu dữ liệu sai được tổng hợp vào bản đồ dịch tễ mà không có cơ chế lọc và xác minh, bản đồ đó có thể gây ra những cảnh báo nhầm và làm mất lòng tin của người dùng. Cần thiết kế một pipeline xác thực chất lượng dữ liệu — có thể kết hợp giữa ngưỡng tin cậy của mô hình, cơ chế đối chiếu chéo giữa nhiều báo cáo cùng vùng, và vòng phản hồi từ chuyên gia nông nghiệp. Dù còn nhiều thách thức, đây vẫn là hướng phát triển xứng đáng đầu tư về dài hạn, bởi nó có thể chuyển hóa hệ thống từ một công cụ chẩn đoán cá nhân đơn thuần thành một nền tảng dữ liệu nông nghiệp cộng đồng với giá trị lan tỏa vượt ra ngoài từng người dùng đơn lẻ.
\clearpage